\chapter{Overview}

\ChapterMeta{
This chapter introduces \yolozu{}'s interface-contract-first evaluation paradigm (predictions-first artifacts for continual learning and test-time adaptation), delineates the primary repository entry points, and clarifies the distinction between its two Command Line Interface (CLI) surfaces.
}{

It is designed to expedite workflow selection and eliminate ambiguity regarding command execution and artifact management.
}{
No prerequisites are required beyond reading. Rely on the \cmd{--help} flag for the specific CLI in use, and treat JSON artifacts as the definitive, stable interface.
}
\ChapterAudience{Start here if you are new to YOLOZU and need the shortest route to the right workflow and CLI surface.}


\section{What \yolozu{} Is}
\yolozu{} is a framework-agnostic evaluation toolkit for vision models,
designed to address \textbf{catastrophic forgetting} and \textbf{domain shift}
through reproducible continual learning and test-time adaptation.

\paragraph{The shortest human summary.}
If you only remember one sentence from this chapter, let it be this:
\begin{quote}
\yolozu{} is the part of the stack that keeps model outputs comparable and reproducible even when the underlying training or inference stack changes.
\end{quote}

That is why the rest of the manual keeps coming back to the same habits:
\begin{itemize}
  \item save artifacts,
  \item validate them early,
  \item compare them under pinned settings,
  \item do not let backend differences silently turn into research drift.
\end{itemize}

\paragraph{\yolozu{} at a glance.}
\begin{itemize}
  \item Framework-agnostic evaluation toolkit for vision models.
  \item Focuses on catastrophic forgetting and domain shift.
  \item Supports continual learning and test-time adaptation/training.
  \item Predictions-first interface contract (JSON artifacts).
  \item Multi-task: detection, segmentation, keypoints, monocular depth, and 6DoF pose.
  \item Deployment + interoperability: ONNX export and runtime backends.
  \item AI-first development: stable CLIs, machine-readable manifests, and regression-testable artifacts.
\end{itemize}

\paragraph{Enterprise-friendly repository posture.}
\yolozu{} keeps the shipped repository code under an Apache-2.0 policy,
does not ship built-in telemetry or phone-home style usage collection in its
tooling, and relies on explicit quality-control mechanisms such as tests,
manifests, CI gates, and provenance/reporting helpers.
This improves enterprise deployability, but it does not remove the separate
review boundary for datasets, model weights, vendor runtimes, or hosted
deployment environments.

\textbf{\yolozu{} exists to measure, compare, and enable mitigation of catastrophic forgetting
and inference-time adaptation in a reproducible way.}

\yolozu{} treats \textbf{predictions---not models---as the stable interface contract},
making continual learning and test-time training comparable,
restartable, and CI-friendly across frameworks and runtimes.

\yolozu{} supports evaluation workflows for object detection, segmentation,
keypoint estimation, \textbf{monocular depth estimation}, and \textbf{6DoF pose estimation},
while keeping training implementations optional and decoupled.

\yolozu{} supports ONNX export and deployment across PyTorch, ONNX Runtime,
and TensorRT, with inference modules in C++ and Rust.

\yolozu{} is built for interface-contract-first, AI-first development:
every experiment produces versioned artifacts that can be
compared and regression-tested automatically in CI.


\section{AI-First Design Intent}
This repository is architected as an \textbf{AI-first} codebase.
In practical terms, this implies:
\begin{itemize}
  \item AI agents must be capable of \textbf{utilizing} the codebase safely via stable command surfaces and well-defined interface contracts.
  \item AI agents must be able to \textbf{extend} workflows (e.g., introducing new experiments, scripts, or reports) without violating existing interfaces.
  \item JSON schemas and machine-readable registries are treated as first-class interfaces, rather than optional metadata.
\end{itemize}

Operational guidance for both human operators and agentic workflows dictates that artifacts must be treated as first-class entities, interface contracts must be validated early in the pipeline, and evaluation protocols must remain strictly pinned.

High-level objectives include:
\begin{itemize}
  \item Generating stable, versioned JSON artifacts (encompassing predictions, evaluation reports, and run metadata).
  \item Ensuring CPU-minimum compatibility for development and testing; GPU acceleration remains optional for training and inference.
  \item Facilitating backend-agnostic workflows (supporting PyTorch, ONNX Runtime, TensorRT, and custom C++/Rust implementations).
  \item Streamlining migration by accepting prevalent dataset ecosystems (e.g., YOLO-style data.yaml formats, COCO JSON) via lightweight wrapper artifacts, thereby preserving the integrity of original datasets.
  \item Enforcing a production-oriented run interface contract (Run Contract) for training operations, which standardizes artifact paths, resumption capabilities, and model exports.
  \item Maintaining a strict Apache-2.0 operational posture, ensuring no non-Apache dependencies are present in shipped code paths.
  \item Avoiding built-in telemetry in shipped tooling; quality control is performed through explicit, reviewable checks instead.
\end{itemize}


\section{Navigating the Repository}
The primary entry points for navigating the project are:
\begin{itemize}
  \item \path{README.md}: Provides a concise overview and essential canonical commands.
  \item \path{docs/README.md}: Serves as the documentation index, organized by \textit{task objective}.
  \item \path{docs/yolozu_spec.md}: Contains the comprehensive feature summary and technical specification.
  \item \path{tools/}: Houses power-user scripts and the repository's wrapper CLI.
\end{itemize}

\paragraph{Module path policy.}
Canonical Python modules are organized in categorized packages such as
\path{yolozu/core}, \path{yolozu/datasets}, \path{yolozu/eval},
\path{yolozu/inference}, \path{yolozu/predictions}, \path{yolozu/training}, and
\path{yolozu/geometry}. New code should import these canonical paths directly.
Legacy top-level imports are preserved through package-level aliasing for compatibility.

\section{How to use this manual}
Different readers usually arrive with different first questions. A quick routing guide helps:
\begin{itemize}
  \item \textbf{``I want to run something today.''} Start with Chapters~2, 4, and 5.
  \item \textbf{``I already have predictions and only need evaluation.''} Start with Chapters~3 and 5.
  \item \textbf{``I care about forgetting, TTT, or research ablations.''} Start with Chapters~13, 14, and 15.
  \item \textbf{``I care about geometry, depth, or pose refinement.''} Start with Chapters~10 and 16.
  \item \textbf{``I am modifying tools or documentation.''} Start with Chapters~12 and 18.
\end{itemize}

\paragraph{A good first reading order for most users.}
\begin{enumerate}
  \item Chapter~1: understand what kind of repository this is.
  \item Chapter~2: get to a working environment.
  \item Chapter~5: choose the right evaluate/export workflow.
  \item Only then branch into the more specialized chapters.
\end{enumerate}


\section{Dual CLI Surfaces: Pip vs. Repository}
\yolozu{} exposes two distinct command surfaces to accommodate different user needs:
\begin{itemize}
  \item \textbf{Pip CLI}: Invoked via \cmd{yolozu ...}, this is the install-safe, user-facing command interface.
  \item \textbf{Repository CLI}: Invoked via \cmd{python3 tools/yolozu.py ...}, this interface is tailored for research, evaluation, and advanced workflows.
\end{itemize}

Both CLIs share underlying concepts and data interface contracts. This manual references whichever CLI is most appropriate for the specific workflow under discussion.

\paragraph{Rule of thumb.}
If you are unsure which CLI to start with:
\begin{itemize}
  \item use \cmd{yolozu ...} when acting as a user,
  \item use \cmd{python3 tools/yolozu.py ...} when working inside the repository and you want the research surface.
\end{itemize}
